% =====================================================================
%  Why Graph-Structured Systems Engaged by Finite Agents
%  Must Admit Operational Structure: A Conditional Derivation
%
%  Author: Kundai Farai Sachikonye
%          Technical University of Munich
% =====================================================================

\documentclass[11pt,a4paper]{article}

\usepackage[utf8]{inputenc}
\usepackage[T1]{fontenc}
\usepackage{lmodern}
\usepackage[a4paper,margin=2.4cm]{geometry}
\usepackage{amsmath,amssymb,amsthm,mathtools}
\usepackage{bm}
\usepackage{mathrsfs}
\usepackage{enumitem}
\usepackage{booktabs}
\usepackage{microtype}
\usepackage[round,sort&compress]{natbib}
\usepackage[colorlinks=true,linkcolor=blue!45!black,
            citecolor=blue!45!black,urlcolor=blue!45!black]{hyperref}
\usepackage{cleveref}

% ---- Theorem environments ----
\newtheoremstyle{thmstyle}{\topsep}{\topsep}{\itshape}{}{\bfseries}{.}{.5em}{}
\theoremstyle{thmstyle}
\newtheorem{theorem}{Theorem}[section]
\newtheorem{lemma}[theorem]{Lemma}
\newtheorem{proposition}[theorem]{Proposition}
\newtheorem{corollary}[theorem]{Corollary}
\theoremstyle{definition}
\newtheorem{definition}[theorem]{Definition}
\newtheorem{axiom}{Premise}
\newtheorem{principle}[theorem]{Principle}
\theoremstyle{remark}
\newtheorem{remark}[theorem]{Remark}

\crefname{axiom}{Premise}{Premises}
\Crefname{axiom}{Premise}{Premises}

% ---- Notation ----
\newcommand{\R}{\mathbb{R}}
\newcommand{\N}{\mathbb{N}}
\newcommand{\Z}{\mathbb{Z}}
\newcommand{\Zp}{\mathbb{Z}_{\geq 0}}
\newcommand{\Graph}{\mathcal{G}}
\newcommand{\Vrt}{\mathsf{V}}
\newcommand{\Edg}{\mathsf{E}}
\newcommand{\Agent}{\mathcal{A}}
\newcommand{\Op}{\mathsf{Op}}
\newcommand{\Res}{\rho}
\newcommand{\floor}{\beta}
\newcommand{\OpStr}{\mathfrak{O}}

\title{\textbf{Why Graph-Structured Systems Engaged by Finite Agents\\
Must Admit Operational Structure}\\[4pt]
\large A Conditional Derivation}

\author{%
Kundai Farai Sachikonye\\
\small Technical University of Munich\\
\small \texttt{kundai.sachikonye@wzw.tum.de}}

\date{\today}

% =====================================================================
\begin{document}
\maketitle

\begin{abstract}
\noindent
We give a conditional derivation: granted two premises---that the
engaging agents are finite, and that the engaged system is a finite
graph---we show that the system necessarily admits an
\emph{operational structure} in a precise sense. The argument is
constructive. We define a finite graph-structured system and a finite
agent, derive that any distinguishability relation an agent can
realise has a strictly positive floor on edge measure, prove that
this floor forces a residue on every agent--graph interaction, and
show that the realised acts form a small category whose hom-sets
carry a preorder by residue. This category, presented with its
residue preorder and a budget-bounded notion of realisability, is
what we call the system's operational structure. The structure is
not chosen; it is the unique closure of the agent's realisable
interactions under the constraints imposed by finiteness on both
sides. We classify operations into four canonical types---algebraic,
dynamical, computational, and internal---and show every realisable
operation is a combination of these. The derivation makes no
metaphysical commitments and applies wherever the two premises
hold; the framework is silent outside that scope. We close by
stating the narrow falsifiability of the main theorem: it can be
refuted only by exhibiting a finite agent realising
distinguishability with zero floor on a finite graph, which we prove
is structurally impossible. The paper is self-contained; no result
of the author is cited; the literature references are to standard
mathematical lemmas.
\end{abstract}

\medskip
\noindent\textbf{Keywords:} finite agents; finite graphs; operational
structure; distinguishability floor; residue propagation; small
category; partial composition; residue preorder; conditional
derivation.

\tableofcontents

% =====================================================================
\section{Introduction}
\label{sec:intro}
% =====================================================================

\subsection{The conditional question}

The question this paper addresses is short:
\begin{quote}\itshape
On the condition that the engaging agents are finite and the engaged
system is a finite graph, what structure of operations must the
system admit?
\end{quote}
We answer: a specific operational structure---a small category
with a residue preorder, a positive floor, and a budget-bounded
realisability marker---derivable from the two premises alone, with
no further hypotheses.

The result is conditional throughout. We do not assert that all
agents are finite, or that all systems are graphs, or that systems
are fundamentally relational. We assert only that \emph{if} both
conditions hold, the operational structure of
\cref{thm:opstructure} is forced. The framework is silent outside
that scope: where either premise fails, the derivation does not
apply and makes no prediction.

\subsection{Why this matters}

Engineering practice routinely treats operational structures---
function composition, command pipelines, state-machine transitions,
query algebras---as design choices to be selected for convenience.
The argument here shows that for the dominant case (finite agents
engaging finite graphs, which covers essentially all of practical
software, embedded systems, sensor networks, knowledge graphs, and
finite databases) the operational structure is \emph{not} a design
choice but a derivable necessity. The class of admissible structures
is forced; only the labelling within the class is free.

This converts a large family of engineering decisions from
\emph{which structure to invent} into \emph{which structure has the
two premises already determined}. The latter is a discovery task
with a definite answer; the former is a design task with no
objective ground.

\subsection{Method}

The derivation is constructive. We fix the two premises in
\cref{sec:premises}, derive three consequences in
\cref{sec:consequences} (distinguishability floor, residue
propagation, small-category structure under partial composition),
prove the main theorem in \cref{sec:main}, classify the resulting
operations in \cref{sec:classification}, and locate the
falsifiability of the argument in \cref{sec:falsifiability}. Every
step is a finite combinatorial, measure-theoretic, or
category-theoretic claim with a self-contained proof; no result of
the author is cited.

\subsection{Relation to existing frameworks}

The argument touches finite graph theory \citep{diestel2017,
bollobas1998}, finite-state automata \citep{hopcroft2007}, measure
theory of finite partitions \citep{halmos1950}, category theory and
partial composition \citep{maclane1998}, semigroup and monoid theory
\citep{howie1995}, and order theory on finite lattices
\citep{daveypriestley2002}. We use only the textbook content of
each field. The novel contribution is the closure argument
(\cref{thm:opstructure}) that forces a small-category operational
structure with floor, residue, and preorder from the conjunction of
the two premises; the constituent pieces are standard, but their
conjunction has not been packaged as a forced structural result, to
the author's knowledge.

% =====================================================================
\section{Premises}
\label{sec:premises}
% =====================================================================

We fix two premises. Both are operational conditions on the agent
and the system; neither is a metaphysical hypothesis about the
world.

\subsection{Finite graph-structured system}

\begin{definition}[Finite graph-structured system]
\label{def:graph}
A \emph{finite graph-structured system} is a triple $\Graph =
(\Vrt, \Edg, w)$ where $\Vrt$ is a finite non-empty set of
\emph{vertices}, $\Edg \subseteq \Vrt \times \Vrt$ is a finite set
of \emph{edges}, and $w: \Edg \to \R_{\geq 0}$ is an \emph{edge
measure} assigning to each edge a non-negative real weight. We do
not require edges to be undirected or distinct; multigraphs are
admitted.
\end{definition}

\begin{axiom}[Finite graph hypothesis]
\label{ax:graph}
The engaged system $\Graph = (\Vrt, \Edg, w)$ is a finite
graph-structured system: $|\Vrt| < \infty$, $|\Edg| < \infty$, and
$w$ takes finite values throughout.
\end{axiom}

The graph need not be connected, simple, or planar. It need not be
metric, in the sense that $w$ need not satisfy the triangle
inequality. We require only that it has a finite number of
vertices, a finite number of edges, and a finite-valued non-negative
measure on its edges.

\begin{remark}[Why graphs]
\label{rem:why_graphs}
We pick graphs as the primitive structure because they are the
minimal combinatorial encoding of relational data: a set of entities
together with the relations between them. Any system describable as
``entities standing in relations'' is encodable as a graph, possibly
a multigraph with labelled edges. Trees, finite-state machines,
knowledge bases, relational databases, dependency graphs, and finite
hypergraphs all reduce to or extend this minimal form. The choice
restricts neither the scope of application nor the generality of the
result.
\end{remark}

\subsection{Finite agent}

\begin{definition}[Distinguishability relation]
\label{def:distrel}
A \emph{distinguishability relation} on $\Graph$ is a binary
relation $\sim \subseteq \Vrt \times \Vrt$ such that for each pair
$(u, v)$, the agent has either established $u \sim v$ (the pair is
treated as a single entity), or $u \not\sim v$ (the pair is treated
as distinct entities), but not both.
\end{definition}

\begin{definition}[Finite agent]
\label{def:agent}
A \emph{finite agent} $\Agent$ on $\Graph$ is specified by:
\begin{enumerate}[label=(F\arabic*),leftmargin=2.6em,nosep]
\item\label{F:fin_state} A finite internal state space $S_\Agent$
with $|S_\Agent| < \infty$, encoding all distinctions the agent can
hold at once.
\item\label{F:fin_descr} A finite descriptive scheme: each
distinguishability decision is recorded by an injection from the
$\sim$-classes into $S_\Agent$, so the number of distinct
equivalence classes the agent can simultaneously hold is at most
$|S_\Agent|$.
\item\label{F:fin_cost} A finite cost budget: each act of
distinguishing a pair $(u,v)$ as $u \not\sim v$ has a cost
$c(u,v) > 0$ bounded below by a positive constant $c_0$, and the
agent's total expenditure over its lifetime is bounded by a finite
budget $C_{\Agent} < \infty$.
\end{enumerate}
\end{definition}

\begin{axiom}[Finite agent hypothesis]
\label{ax:agent}
Any agent engaging $\Graph$ is a finite agent in the sense of
\cref{def:agent}: it has finite internal state, finite descriptive
scheme, and finite cost budget.
\end{axiom}

We make no further assumptions about the agent's architecture: it
may be a human, a machine, a sensor network, a software process, a
finite collective. The three finiteness conditions of
\cref{def:agent} are the substantive content of \cref{ax:agent}.

\subsection{What the premises do not say}

\Cref{ax:graph,ax:agent} are sharp where they need to be sharp and
silent elsewhere. We list three things the premises explicitly do
not say.

\begin{enumerate}[nosep,leftmargin=*]
\item They do not say the world is finite. Reality at large is
permitted to be infinite; the premises restrict only the system the
agent engages and the agent itself, both of which are finite.
\item They do not say the graph encodes everything about the system.
A graph is a \emph{model} of the system the agent uses; what the
system is in itself, outside the agent's modelling, is not
addressed.
\item They do not say agents must be physical. The argument applies
to any agent satisfying~\ref{F:fin_state}--\ref{F:fin_cost},
whatever its substrate.
\end{enumerate}

\subsection{Scope}

The conjunction of the two premises covers the dominant case of
engineering practice: software systems, embedded controllers,
sensor networks, knowledge bases, finite databases, finite-state
protocols, and finite communication graphs. It does not cover
putative infinite agents (perfect oracles, unbounded computation) or
putative infinite systems (continuous fields engaged directly,
without finite discretisation). Where finite discretisation is
imposed---which is the case in all realised engineering
systems---the premises apply.

% =====================================================================
\section{Consequences of the Two Premises}
\label{sec:consequences}
% =====================================================================

We derive three consequences of the conjunction of \cref{ax:graph}
and \cref{ax:agent}. Each is forced by the premises alone; no
additional hypothesis is invoked.

\subsection{The distinguishability floor}

\begin{theorem}[Floor on edge measure for realisable
distinguishability]
\label{thm:floor}
Let $\Graph$ satisfy \cref{ax:graph} and $\Agent$ satisfy
\cref{ax:agent}. Define the graph-level floor
\[
  \floor := \min\{w(e) : e \in \Edg,\; w(e) > 0\}.
\]
This is a property of $\Graph$ alone, well-defined whenever
$\Graph$ contains at least one positive-measure edge. Then
$\floor > 0$, and for every pair $(u,v) \in \Vrt \times \Vrt$ that
$\Agent$ can realise as $u \not\sim v$, the agent's
distinguishability path through $\Graph$ contains an edge of measure
$\geq \floor$.
\end{theorem}

\begin{proof}
\emph{Positivity of $\floor$.} The set
$W^{+} := \{w(e) : e \in \Edg,\; w(e) > 0\}$ is a finite (by
$|\Edg| < \infty$ in \cref{ax:graph}) subset of $\R_{> 0}$. If
$W^{+}$ is non-empty, $\floor = \min W^{+} > 0$. If $W^{+}$ is
empty, every edge has measure $0$, no path can have positive
total measure, and \cref{F:fin_descr} prevents the agent from
recording any non-trivial distinguishability (a null path gives the
agent no measure differential to record); in this degenerate case
no realisable distinguishability acts exist and the claim holds
vacuously.

\emph{Realised paths contain a $\geq \floor$ edge.} Assume the
non-degenerate case, so $\floor > 0$. Let $(u,v)$ be realised by
$\Agent$ as $u \not\sim v$. By \cref{F:fin_descr}, the agent
records the distinction by some path
$P_\Agent(u,v) = (e_1, e_2, \ldots, e_k)$ in $\Graph$. If every
edge of $P_\Agent(u,v)$ had measure $0$, the path would have total
measure $0$, the agent's descriptive scheme would have no measure
differential to encode the distinction, and the distinguishability
would not be realised---contradicting the assumption. Hence at
least one edge $e_i$ of $P_\Agent(u,v)$ has $w(e_i) > 0$, and
therefore $w(e_i) \geq \floor$ by the definition of $\floor$ as
the minimum positive edge measure. $\qed$
\end{proof}

\begin{remark}[The floor is a graph property, not an agent
property]
\label{rem:floor_graph_property}
$\floor$ depends only on $\Graph$, not on $\Agent$. Different
finite agents engaging the same graph share the same floor.
$\Agent$ enters only in the second clause, where realisability
guarantees that at least one edge above the floor lies on every
realised distinguishability path.
\end{remark}

\begin{remark}[Floor is a finiteness floor]
\label{rem:floor_is_finiteness}
The floor $\floor$ is not a metaphysical lower bound on what edges
``can be.'' It is the lower bound forced by the finiteness of
$\Edg$ and the requirement that distinguishability be realisable.
With a finite edge set having positive-measure edges, the minimum
positive measure is itself positive. This is a combinatorial
observation, not a substantive claim about the structure of
reality.
\end{remark}

\subsection{Residue: every distinguishability act deposits content}

\begin{definition}[Residue]
\label{def:residue}
The \emph{residue} $\Res(u,v)$ of a distinguishability act
$u \not\sim v$ realised by $\Agent$ is the sum of edge measures
along the path the agent uses:
\[
  \Res(u,v) = \sum_{e \in \mathrm{path}_\Agent(u,v)} w(e).
\]
By \cref{thm:floor}, $\Res(u,v) \geq \floor > 0$ for every
realisable distinguishability act.
\end{definition}

\begin{theorem}[Residues are additive under composition]
\label{thm:residue_additive}
Let $u \not\sim v$ and $v \not\sim w$ be two distinguishability
acts realised by $\Agent$. The composed act $u \not\sim w$ (if
realisable) has residue:
\[
  \Res(u, w) \leq \Res(u, v) + \Res(v, w).
\]
Equality holds when the path the agent uses for the composed act is
the concatenation of the individual paths.
\end{theorem}

\begin{proof}
The residue of the composed act is the sum of edge measures along
the agent's path from $u$ to $w$. By the triangle inequality for
sums of non-negative numbers, the shortest such path has total
measure no greater than the sum of any concatenated route. If the
agent uses the concatenation $u \to v \to w$, equality holds;
otherwise the agent has found a shorter path and the inequality is
strict. $\qed$
\end{proof}

\begin{corollary}[Residue accumulates and does not refund]
\label{cor:res_accumulate}
For any sequence of distinguishability acts
$u_0 \not\sim u_1 \not\sim u_2 \not\sim \cdots \not\sim u_n$, the
total residue deposited satisfies:
\[
  \sum_{i=0}^{n-1} \Res(u_i, u_{i+1}) \geq n\floor > 0.
\]
The cumulative residue strictly grows with the number of acts and
is bounded below by $n\floor$. No act reduces the cumulative total;
the sum is monotone non-decreasing in $n$.
\end{corollary}

\begin{proof}
By \cref{thm:floor}, each $\Res(u_i, u_{i+1}) \geq \floor$. Summing
$n$ non-negative terms each $\geq \floor$ gives $\geq n\floor$.
Monotonicity in $n$ is immediate from the non-negativity of each
term. $\qed$
\end{proof}

\subsection{Closure: distinguishability acts form a small category}

Composition of distinguishability acts is intrinsically
\emph{partial}: $(u,v)$ and $(v,w)$ both being realised by $\Agent$
does not imply that $(u,w)$ is realised. Realisability depends on
the agent's path-choices and on the budget $C_\Agent$. The
finite-budget premise \cref{F:fin_cost} actively limits which
compositions are realised, so a total composition law cannot be
assumed.

The honest object is therefore a \emph{small category}, not a
monoid. In a small category, composition is partial in exactly
this controlled sense: a morphism $g: u \to v$ composes with a
morphism $f: v \to w$ to give $f \circ g: u \to w$, but morphisms
with mismatched source/target do not compose. We construct the
category and verify the axioms.

\begin{definition}[The distinguishability category]
\label{def:dcat}
The \emph{distinguishability category}
$\mathbf{C}_{\Agent,\Graph}$ has:
\begin{itemize}[nosep,leftmargin=*]
\item \emph{Objects}: the vertices $\Vrt$ of $\Graph$.
\item \emph{Morphisms}: a morphism $u \to v$ is either the
identity $\mathrm{id}_u$ (with residue $0$), or a realised
distinguishability act $(u,v) \in \mathcal{D}_\Agent$ together with
the agent's path $\mathrm{path}_\Agent(u,v)$ used to realise it.
\item \emph{Composition}: for morphisms $g: u \to v$ and
$h: v \to w$, the composite $h \circ g: u \to w$ is defined iff
$\Agent$ realises a distinguishability act $(u,w)$ whose path is
recordable from the agent's record of $g$ and $h$. When defined,
$h \circ g = (u, w, \mathrm{path}_\Agent(u,w))$.
\item \emph{Identities}: for each $u \in \Vrt$, the identity
$\mathrm{id}_u : u \to u$ is the trivial act with empty path and
residue $0$.
\end{itemize}
\end{definition}

\begin{lemma}[Identity laws]
\label{lem:identity}
For every realised morphism $g: u \to v$,
$g \circ \mathrm{id}_u = g$ and $\mathrm{id}_v \circ g = g$.
\end{lemma}

\begin{proof}
The composites $g \circ \mathrm{id}_u$ and $\mathrm{id}_v \circ g$
have source $u$, target $v$, residue $\Res(g) + 0 = \Res(g)$, and
path $\mathrm{path}_\Agent(u,v)$. Both are the realised morphism
$g$. The identity composites are realisable because they require
no agent effort beyond what $g$ itself already required. $\qed$
\end{proof}

\begin{lemma}[Associativity where defined]
\label{lem:assoc}
For realised morphisms $g_1: u \to v$, $g_2: v \to w$,
$g_3: w \to x$ such that both $(g_3 \circ g_2) \circ g_1$ and
$g_3 \circ (g_2 \circ g_1)$ are defined,
\[
  (g_3 \circ g_2) \circ g_1 \;=\; g_3 \circ (g_2 \circ g_1).
\]
\end{lemma}

\begin{proof}
Both composites have source $u$ and target $x$. By
\cref{F:fin_descr}, the agent records a morphism by its
endpoints and its realised path; if both composites are defined,
the agent has recorded a realised morphism $u \to x$, and the
record is unique up to the agent's path-choice. The agent's
path-choice for the composite is a function only of the endpoints
and the available realised sub-morphisms, which both bracketings
present identically. Hence the two composites are the same
morphism. $\qed$
\end{proof}

\begin{theorem}[Realised acts form a small category]
\label{thm:category}
Under \cref{ax:graph} and \cref{ax:agent},
$\mathbf{C}_{\Agent,\Graph}$ is a small category: objects form a
finite set ($|\Vrt| < \infty$), morphisms form a finite set
(bounded by $|\Vrt| + |\mathcal{D}_\Agent| \leq |\Vrt| +
|S_\Agent|^2 < \infty$), composition is partial and associative
where defined (\cref{lem:assoc}), and identities exist
(\cref{lem:identity}).
\end{theorem}

\begin{proof}
Finiteness of objects and morphisms is immediate from the premises.
\cref{lem:identity,lem:assoc} provide the identity and associativity
axioms in the partial-composition setting. $\qed$
\end{proof}

\begin{remark}[Why a category rather than a monoid]
\label{rem:why_category}
A monoid requires composition to be total. In our setting,
realisability depends on the agent's finite budget and on whether
the agent's recording of one act can be extended to a recording of
a composite act. Both fail for some triples. The category framework
is the standard mathematical home for partial composition with
identities and associativity \citep{maclane1998}; it captures
exactly what the premises support, without overclaiming totality.
A small category in which every pair of composable morphisms does
compose is a monoid; the framework reduces to a monoid in that
special case (no budget exhaustion, full closure), but the general
case is a category.
\end{remark}

\begin{remark}[Free monoid envelope]
\label{rem:free_monoid}
The path category over $\mathbf{C}_{\Agent,\Graph}$---formal
sequences of composable morphisms, taken modulo the realised
composites---is a free monoid on the morphisms of
$\mathbf{C}_{\Agent,\Graph}$, modulo the relations imposed by
realisation. This is a derived construction, useful for some
applications, but not the primitive structure forced by the
premises. The primitive structure is the small category.
\end{remark}

% --- Hom-set sizes ---

\begin{proposition}[Hom-set bound]
\label{prop:hom_bound}
For any $u, v \in \Vrt$, the hom-set
$\mathrm{Hom}_{\mathbf{C}_{\Agent,\Graph}}(u, v)$ contains at most
one realised non-identity morphism plus (if $u = v$) the identity,
totalling at most two elements.
\end{proposition}

\begin{proof}
A realised distinguishability act $(u, v)$ is the assertion
$u \not\sim v$; the agent records this assertion at most once in
its descriptive scheme (it is either established or not). If the
agent's path-choice changes during a run, the agent's record is
updated, but the morphism remains the unique realised act between
$u$ and $v$. The identity on $u$ exists when $u = v$. $\qed$
\end{proof}

% =====================================================================
\section{The Operational Structure Theorem}
\label{sec:main}
% =====================================================================

We now state and prove the main result: a finite agent engaging a
finite graph cannot avoid admitting an operational structure in the
following precise sense.

\subsection{Operational structure}

The previous section established three things: every realised
distinguishability act has residue $\geq \floor > 0$
(\cref{thm:floor}); residues compose subadditively
(\cref{thm:residue_additive}); the realised acts form a small
category $\mathbf{C}_{\Agent,\Graph}$ (\cref{thm:category}). The
operational structure packages these into a single object.

\begin{definition}[Residue equivalence and residue quotient]
\label{def:residue_equiv}
Two morphisms $f, g$ of $\mathbf{C}_{\Agent,\Graph}$ with
$\Res(f) = \Res(g)$ are \emph{residue-equivalent}, written
$f \equiv_\Res g$. The relation $\equiv_\Res$ is an equivalence on
each hom-set; the \emph{residue quotient}
$\widetilde{\mathbf{C}}_{\Agent,\Graph}$ identifies
residue-equivalent morphisms.
\end{definition}

\begin{definition}[Operational structure]
\label{def:opstr}
The \emph{operational structure} on a finite graph $\Graph$
relative to a finite agent $\Agent$ is the quintuple
\[
  \OpStr_\Agent \;=\;
  (\mathbf{C}_{\Agent,\Graph},\; \Res,\; \preceq,\; \floor,\;
  \mathrm{Realisable})
\]
where:
\begin{enumerate}[label=(O\arabic*),leftmargin=2.6em,nosep]
\item\label{O:cat} $\mathbf{C}_{\Agent,\Graph}$ is the small
category of \cref{def:dcat}.
\item\label{O:residue} $\Res: \mathrm{Mor}(\mathbf{C}_{\Agent,
\Graph}) \to \R_{\geq 0}$ is the residue function of
\cref{def:residue}, subadditive under composition
(\cref{thm:residue_additive}).
\item\label{O:preorder} $\preceq$ is the residue \emph{preorder}:
$f \preceq g$ iff $\Res(f) \leq \Res(g)$. On the residue quotient
$\widetilde{\mathbf{C}}_{\Agent,\Graph}$ this descends to a partial
order in the standard sense.
\item\label{O:floor} For every non-identity morphism $f$,
$\Res(f) \geq \floor > 0$, where $\floor$ is the graph-level floor
of \cref{thm:floor}.
\item\label{O:budget} $\mathrm{Realisable}: \mathrm{Mor}(\mathbf{C}_
{\Agent,\Graph}) \to \{0,1\}$ marks each morphism as realised or
not, subject to the budget constraint
$\sum_{f\text{ realised}} c(f) \leq C_\Agent$.
\end{enumerate}
\end{definition}

\begin{remark}[Why these five components]
\label{rem:five}
The category (\ref{O:cat}) captures composition; the residue
function (\ref{O:residue}) captures cost; the preorder
(\ref{O:preorder}) captures comparison; the floor (\ref{O:floor})
captures the irreducible minimum; the realisability marker
(\ref{O:budget}) captures the budget-constrained partiality. None
of these can be dropped without losing structural content forced by
the premises, and none requires more than the premises supply.
Together they describe what the agent can do, at what cost, in what
order, with what minimum, within what budget.
\end{remark}

\begin{theorem}[Operational Structure Theorem]
\label{thm:opstructure}
Under \cref{ax:graph} and \cref{ax:agent}, every distinguishability
relation $\sim$ that the agent realises on $\Graph$ determines a
unique operational structure $\OpStr_\Agent$. The structure is
forced by the conjunction of the two premises: each of its five
components is determined by $(\Agent, \Graph)$ together with the
agent's path-choices, and no component can be denied without
denying one of the premises.
\end{theorem}

\begin{proof}
We verify each clause and the uniqueness and forcing claims.

\emph{Category~\ref{O:cat}.} By \cref{thm:category},
$\mathbf{C}_{\Agent,\Graph}$ is a small category with object set
$\Vrt$, identity morphisms on each object, partial associative
composition. The morphisms include the identities and the
agent's realised distinguishability acts together with their paths.

\emph{Residue~\ref{O:residue}.} The residue function of
\cref{def:residue} assigns to each morphism the sum of edge
measures along its path. Subadditivity under composition is
\cref{thm:residue_additive}.

\emph{Preorder~\ref{O:preorder}.} The relation $f \preceq g$ iff
$\Res(f) \leq \Res(g)$ is reflexive ($\Res(f) \leq \Res(f)$) and
transitive ($\Res(f) \leq \Res(g) \leq \Res(h)$ implies
$\Res(f) \leq \Res(h)$). It is not antisymmetric on
$\mathbf{C}_{\Agent,\Graph}$ itself: distinct morphisms can have
equal residue, in which case both $f \preceq g$ and $g \preceq f$
hold without $f = g$. The preorder descends to a partial order on
the residue quotient
$\widetilde{\mathbf{C}}_{\Agent,\Graph} = \mathbf{C}_{\Agent,
\Graph}/\!\equiv_\Res$, where antisymmetry is recovered by
construction.

\emph{Floor~\ref{O:floor}.} By \cref{thm:floor}, the graph-level
floor $\floor = \min\{w(e) : e \in \Edg, w(e) > 0\}$ is a property
of $\Graph$ alone, and every realised non-identity morphism has
residue $\geq \floor$. Identity morphisms have residue $0$, so the
floor applies only to non-identity morphisms.

\emph{Realisability~\ref{O:budget}.} The realisability marker is
determined by the agent's actual run on $\Graph$, subject to
\cref{F:fin_cost}. Compositions of realised morphisms are
themselves realised only when the agent can record the composite
within the remaining budget; this is the partial-composition
phenomenon noted in \cref{rem:why_category}.

\emph{Uniqueness.} Each of the five components is determined by
$(\Agent, \Graph)$ together with the agent's path-choices: the
category is fixed by the realised acts; the residue function is
fixed by $w$; the preorder is fixed by residue comparison; the
floor is fixed by $\min\{w(e) : w(e) > 0\}$; the realisability
marker is fixed by the agent's run subject to the budget. Hence
$\OpStr_\Agent$ is unique given $(\Agent, \Graph)$ and the path
record.

\emph{Forcing.} The five clauses follow from the premises alone:
rejecting~\ref{O:cat} requires either infinite objects (denial of
\cref{ax:graph}) or non-associative composition (denial of
\cref{lem:assoc}, which uses only the agent's deterministic
recording, \cref{F:fin_descr});
rejecting~\ref{O:residue} requires negative or undefined edge
measures (denial of \cref{def:graph});
rejecting~\ref{O:preorder} requires denying the standard order on
$\R_{\geq 0}$;
rejecting~\ref{O:floor} requires zero-floor distinguishability acts
(denial of \cref{thm:floor}, hence of finiteness of $\Edg$);
rejecting~\ref{O:budget} requires unbounded budget (denial of
\cref{F:fin_cost}). The structure is forced. $\qed$
\end{proof}

\begin{remark}[Strengthening to a partial monoid]
\label{rem:partial_monoid}
If the agent's budget $C_\Agent$ is large enough that composition
of realised morphisms is always realised (no budget exhaustion on
composites), the category becomes a small category in which every
pair of composable morphisms composes---a one-object case of this
being a monoid. The general statement above does not assume this;
it admits the budget-exhausted case where some composites are not
realised. The strengthened result is a corollary under the
additional hypothesis $C_\Agent = \infty$ (or sufficient for
closure), which we do not adopt by default.
\end{remark}

\begin{remark}[Path-choice dependence]
\label{rem:path_choice}
$\OpStr_\Agent$ depends on the agent's path-choices because two
different paths between the same vertices can give different
residues. This is genuine: an agent choosing a longer path pays
more residue. We do not normalise to the shortest path; the
operational structure records what the agent actually did, not
what it could have done. To compare agents at the same graph, one
fixes a path-choice rule (e.g.~shortest-path policy) and the
remaining freedom collapses.
\end{remark}

\subsection{Consequences of the theorem}

\begin{corollary}[Operations are determined up to labelling and
path-choice]
\label{cor:not_chosen}
The category $\mathbf{C}_{\Agent,\Graph}$, the residue function
$\Res$, the preorder $\preceq$, the floor $\floor$, and the
realisability marker are determined by $(\Agent, \Graph)$ together
with the agent's path-choices. Choosing an operational structure
for a finite agent on a finite graph is choosing labels for
elements of a structure that already exists and, where the agent
has a choice of path between vertices, choosing the path; it is
not creating the structure.
\end{corollary}

\begin{proof}
Immediate from the uniqueness clause of \cref{thm:opstructure}.
The freedom is two-fold: labelling of objects and morphisms
(arbitrary names), and path-choice between vertices when multiple
paths exist (\cref{rem:path_choice}). The structure underneath is
fixed. $\qed$
\end{proof}

\begin{corollary}[Finite-system engineering is partly a discovery
task]
\label{cor:discovery}
Designing operations for a finite system engaged by finite agents
is partly the task of discovering the structure that
$(\Agent, \Graph)$ determines, and partly the task of fixing
path-choice and labelling conventions. Two designs disagreeing
about which operations to admit are disagreeing about one or more
of: which finite graph the system is modelled as; which finite
agent the program is; which path-choice policy governs route
selection.
\end{corollary}

\begin{proof}
By \cref{cor:not_chosen}, the structure is determined up to
labelling and path-choice. Differences between designs reduce to
differences in $(\Agent, \Graph)$ or in the path-choice policy.
Pure label differences do not affect the structure. $\qed$
\end{proof}

\begin{corollary}[Non-identity morphisms have residue
$\geq \floor$]
\label{cor:no_below}
Every non-identity morphism of $\mathbf{C}_{\Agent,\Graph}$ has
residue $\geq \floor > 0$. The floor is attained: morphisms of
residue exactly $\floor$ exist whenever some single-edge path of
the agent realises a distinction along an edge of weight $\floor$.
The floor cannot be crossed downward by any realisable composition.
\end{corollary}

\begin{proof}
The lower bound is~\ref{O:floor}. Attainment: pick any edge $e$ of
$\Graph$ with $w(e) = \floor$ (such an edge exists by the
definition of $\floor$ in \cref{thm:floor}). If the agent realises
the distinguishability across this single edge, the resulting
morphism has residue $\floor$. Non-crossable downward:
\cref{thm:residue_additive} gives subadditivity, so the residue of
a composite is bounded below by the sum of residues of its
non-identity components, each $\geq \floor$; the sum cannot be
less than $\floor$ unless all components are identities, in which
case the composite is itself the identity. $\qed$
\end{proof}

% =====================================================================
\section{Classification of Operational Structures}
\label{sec:classification}
% =====================================================================

We classify the operations arising in $\OpStr_\Agent$ by the
\emph{locus of their effect}. An operation acts on at most three
loci: the graph $\Graph$, the equivalence structure
$\sim$ on $\Vrt$, and the agent's internal state $S_\Agent$. We
classify by which loci an operation touches. This yields four
canonical types; we prove that these four exhaust the realisable
operations and give the residue lower bound for each.

\subsection{The four loci of effect}

\begin{definition}[Effect tuple]
\label{def:effect_tuple}
The \emph{effect tuple} of a realised operation $f$ is the triple
\[
  \mathrm{eff}(f) \;=\; (\Delta\Graph(f),\; \Delta{\sim}(f),\;
  \Delta S_\Agent(f))
\]
where $\Delta\Graph(f) \in \{0, 1\}$ records whether $f$ modifies
$\Graph$ (vertex/edge set or $w$); $\Delta{\sim}(f) \in \{0, 1\}$
records whether $f$ modifies the agent's equivalence relation on
$\Vrt$; and $\Delta S_\Agent(f) \in \{0, 1\}$ records whether $f$
modifies the agent's internal state $S_\Agent$ beyond what
$\Delta\Graph$ and $\Delta{\sim}$ already entail.
\end{definition}

\begin{remark}[Internal effects beyond the graph]
\label{rem:internal_beyond}
The third component $\Delta S_\Agent$ accounts for internal
bookkeeping not visible in $\Graph$ or $\sim$: counter increments,
timer updates, allocation of internal registers, logging of
realisations, garbage collection of stale records. By
\cref{F:fin_state} the internal state is finite; changes to it are
finite-step. By the finite-budget premise \cref{F:fin_cost} these
changes too consume part of $C_\Agent$, even when no
graph-side or equivalence-side effect occurs.
\end{remark}

\subsection{The four types}

\begin{definition}[Operation type]
\label{def:op_type}
A realised operation $f$ has type:
\begin{enumerate}[label=(T\arabic*),leftmargin=2.6em,nosep]
\item\label{T:alg} \emph{Algebraic} if
$\mathrm{eff}(f) = (0, 1, *)$: $f$ changes only the equivalence
relation on $\Vrt$, not the graph itself.
\item\label{T:dyn} \emph{Dynamical} if
$\mathrm{eff}(f) = (1, *, *)$: $f$ modifies $\Graph$.
\item\label{T:comp} \emph{Computational} if
$\mathrm{eff}(f) = (0, 0, *)$ and $f$ produces a binary output
(yes/no decision) by traversing a finite portion of $\Graph$
without modifying $\Graph$ or $\sim$.
\item\label{T:int} \emph{Internal} if
$\mathrm{eff}(f) = (0, 0, 1)$ and $f$ produces no externally-
visible output: $f$ updates only $S_\Agent$.
\end{enumerate}
The wildcard $*$ in \ref{T:alg} and \ref{T:dyn} indicates that
internal bookkeeping is permitted but not required.
\end{definition}

\paragraph{Algebraic.} Equivalence-class formation, vertex
relabelling under a graph automorphism, partition refinement.
Algebraic operations re-express the agent's distinguishability
relation without editing the graph or emitting output.

\paragraph{Dynamical.} Graph edits: adding or removing vertices,
adding or removing edges, modifying $w$. Dynamical operations
transition $\Graph_t \to \Graph_{t+1}$.

\paragraph{Computational.} Queries: membership tests, reachability
checks, path-existence checks, equivalence tests. A computational
operation returns a binary output by reading the graph.

\paragraph{Internal.} Pure internal-state operations: counter
increments, timer updates, allocation, logging, internal cache
invalidation, garbage collection of agent-side records. An internal
operation costs budget (it consumes finite agent resources) but
deposits no residue on $\Graph$ and produces no externally-visible
output.

\begin{lemma}[Type-specific residue bounds]
\label{lem:type_residues}
\begin{enumerate}[nosep,leftmargin=*]
\item Algebraic operations (\ref{T:alg}): residue $\geq \floor$
when the equivalence change requires reading at least one edge of
weight $\geq \floor$; residue $0$ when the equivalence change is a
relabelling using no edge information (e.g.~renaming with no
distinguishability assertion).
\item Dynamical operations (\ref{T:dyn}): residue $\geq \floor$
when the edit affects at least one edge of weight $\geq \floor$;
residue $0$ when the edit touches only null edges (a vacuous case
that cannot contribute to any distinguishability,
\cref{thm:floor}).
\item Computational operations (\ref{T:comp}): residue $\geq \floor$
when the query reads at least one positive-measure edge; residue
$0$ for queries answered without reading any positive-measure edge
(e.g.~constant queries, queries decidable from $S_\Agent$ alone).
\item Internal operations (\ref{T:int}): residue $0$ on $\Graph$
(by definition no graph access), but cost $\geq c_0 > 0$ on the
agent's budget by \cref{F:fin_cost}.
\end{enumerate}
\end{lemma}

\begin{proof}
Each clause is immediate from the definitions and from
\cref{thm:floor}. The residue bound applies to graph-side
contributions; budget cost is bounded below by $c_0$ for any
recorded operation, including internal ones. $\qed$
\end{proof}

\subsection{Four-type exhaustiveness}

\begin{theorem}[Four-type exhaustiveness]
\label{thm:three_types}
Every realised operation $f \in \mathrm{Mor}(\mathbf{C}_{\Agent,
\Graph})$ has effect tuple in
$\{0,1\}^3 \setminus \{(0,0,0)\}$ and is therefore one of the four
types of \cref{def:op_type}, or a composite whose components fall
into these types. The all-zero tuple $(0,0,0)$ corresponds to the
identity morphism, which is not a non-trivial operation.
\end{theorem}

\begin{proof}
The effect tuple has $2^3 = 8$ possible values. We classify each:
\begin{enumerate}[nosep,leftmargin=*]
\item $(0,0,0)$: no effect anywhere; this is the identity morphism
$\mathrm{id}$.
\item $(0,0,1)$: only internal state changes; type \ref{T:int}
(internal).
\item $(0,1,0)$ and $(0,1,1)$: equivalence relation changes
without graph modification; type \ref{T:alg} (algebraic). The
internal-state component is permitted as bookkeeping
(\cref{rem:internal_beyond}).
\item $(1,0,0)$, $(1,0,1)$, $(1,1,0)$, $(1,1,1)$: graph
modifications, possibly accompanied by equivalence or internal
changes; type \ref{T:dyn} (dynamical), with the wildcard absorbing
the other components.
\end{enumerate}
The eighth case, computational output, is not encoded in
$\mathrm{eff}(f)$ as a fourth component because output by itself
is an internal-state effect (the agent's internal state records the
output). A computational operation that reads $\Graph$ without
modifying it and writes its output to $S_\Agent$ has effect
$(0,0,1)$, which is also the internal type; the distinguishing
feature of computational versus internal is whether the operation
\emph{reads the graph} during its execution. We refine: type
\ref{T:comp} (computational) is the subtype of $(0,0,1)$ in which
the agent traverses at least one edge of $\Graph$ during the
operation; type \ref{T:int} (internal) is the subtype of $(0,0,1)$
in which no graph edge is traversed.

This gives four mutually exclusive types covering all
non-identity effect tuples. Any realised operation is either an
identity, falls into one of the four types, or is a composite
whose components do. $\qed$
\end{proof}

\begin{remark}[The taxonomy is a definition that earns its
``completeness'' from the loci]
\label{rem:taxonomy_honest}
\Cref{thm:three_types} is exhaustive over the loci of effect
specified in \cref{def:effect_tuple}---graph, equivalence,
internal-state. The exhaustiveness is therefore relative to the
three-locus model. If a richer model is adopted (additional loci,
e.g.~external communication channels, environment side-effects),
the classification extends but is not changed in spirit:
loci-effect-typing remains the organising principle, and the
floor bounds continue to apply locus by locus.
\end{remark}

\begin{remark}[Composition across types]
\label{rem:cross_type}
Operations of different types compose where defined: a query
followed by an edit followed by a relabelling is a single
composite morphism in $\mathbf{C}_{\Agent,\Graph}$. The
classification labels the type of each component; composition
operates on the composites. As before, composition is partial
(\cref{rem:why_category}); not every composite of realised
operations is itself realised.
\end{remark}

% =====================================================================
\section{The Forced Structure in Detail}
\label{sec:detail}
% =====================================================================

We restate the main result as a structural recipe, making explicit
what $\OpStr_\Agent$ looks like once $(\Agent, \Graph)$ is fixed.

\subsection{The structure recipe}

Given $(\Agent, \Graph)$ satisfying \cref{ax:agent} and
\cref{ax:graph}, together with a path-choice policy for the agent:
\begin{enumerate}[label=(R\arabic*),leftmargin=2.6em,nosep]
\item\label{R:floor} Compute $\floor = \min \{w(e) : e \in \Edg,\,
w(e) > 0\}$. This is the graph-level floor.
\item\label{R:cat} Enumerate the agent's realised
distinguishability acts on $\Graph$; together with the identity on
each vertex, these are the morphisms of
$\mathbf{C}_{\Agent,\Graph}$.
\item\label{R:residue} For each morphism $f$, compute the residue
$\Res(f)$ as the path-measure sum.
\item\label{R:preorder} Order the morphisms by residue:
$f \preceq g$ iff $\Res(f) \leq \Res(g)$. Quotient by
residue-equivalence $\equiv_\Res$ to obtain a partial order on
$\widetilde{\mathbf{C}}_{\Agent,\Graph}$.
\item\label{R:classify} Classify each morphism by effect tuple
(\cref{def:effect_tuple}) into one of the four types
\ref{T:alg}--\ref{T:int}.
\item\label{R:budget} Mark each morphism as realised or not,
subject to the budget constraint $\sum c(f) \leq C_\Agent$.
\end{enumerate}

\begin{proposition}[The recipe is well-defined and terminates]
\label{prop:recipe}
The recipe \ref{R:floor}--\ref{R:budget} is well-defined and
terminates in finite time for any $(\Agent, \Graph)$ satisfying
the two premises.
\end{proposition}

\begin{proof}
\ref{R:floor} is a minimum over a finite non-empty set of positive
reals: well-defined, computable in $O(|\Edg|)$ time.
\ref{R:cat} enumerates a finite morphism set bounded by
$|\Vrt| + |S_\Agent|^2$, terminating in finite time.
\ref{R:residue} sums finitely many non-negative reals per morphism.
\ref{R:preorder} is a comparison on a finite totally ordered set
($\R_{\geq 0}$); residue-equivalence quotienting is a finite
partition.
\ref{R:classify} is a case distinction over four cases per
morphism.
\ref{R:budget} is a finite sum compared to a finite bound.
Each step terminates; the recipe terminates. $\qed$
\end{proof}

\subsection{What the structure looks like}

\begin{proposition}[Structural shape]
\label{prop:shape}
The operational structure $\OpStr_\Agent$ has the following shape:
\begin{enumerate}[nosep,leftmargin=*]
\item It is a small category (\cref{thm:category}) with finite
object set and finite morphism set, and with partial composition
constrained by the agent's budget.
\item Its non-identity morphisms have residues bounded below by
$\floor > 0$ (\cref{thm:floor}).
\item It carries a residue preorder, with identity morphisms at
residue $0$; on the residue quotient
$\widetilde{\mathbf{C}}_{\Agent,\Graph}$ the preorder descends to a
partial order.
\item Its morphisms decompose into four canonical types
(\cref{thm:three_types}): algebraic, dynamical, computational,
internal.
\end{enumerate}
\end{proposition}

\begin{proof}
Each clause is a restatement of an earlier result. The shape is
the conjunction of those results. $\qed$
\end{proof}

\subsection{Why the structure is operational rather than only
combinatorial}

The structure $\OpStr_\Agent$ is called \emph{operational} because
each morphism represents an act the agent can perform, with a
quantifiable cost (residue) and a quantifiable order (residue
preorder). It differs from a purely combinatorial structure---a
small category considered abstractly---in two ways:
\begin{enumerate}[nosep,leftmargin=*]
\item Each element is associated with a concrete act on the graph,
not merely an abstract symbol.
\item Composition is realised, not merely defined: composing two
acts produces a third act the agent can perform, with predictable
residue deposited on the graph.
\end{enumerate}
The operational character is essential to the result: we are not
classifying the agent's grammar but its realisable acts on a
concrete finite system. The two coincide only because the premises
force the agent's grammar to track its acts; in richer settings
the two diverge.

% =====================================================================
\section{Falsifiability}
\label{sec:falsifiability}
% =====================================================================

\subsection{Where the argument can fail}

The main theorem (\cref{thm:opstructure}) is a conditional claim
under \cref{ax:graph} and \cref{ax:agent}. It can be falsified only
by:
\begin{enumerate}[nosep,leftmargin=*]
\item Showing the two premises jointly imply a contradiction (they
do not, by inspection: a finite graph and a finite agent obviously
coexist; we are reading this on one).
\item Showing some step in the derivation does not follow from the
premises (a mathematical refutation of one of the proved lemmas).
\item Showing the premises do not apply to a particular case (which
is honest: the framework declines to address that case).
\end{enumerate}

The substantive falsification target is (2): exhibiting a finite
agent realising distinguishability with zero floor on a finite
graph. We close this case in the next subsection.

\subsection{Zero-floor distinguishability is structurally impossible}

\begin{theorem}[Zero-floor impossibility]
\label{thm:zero_floor}
Under \cref{ax:graph} and \cref{ax:agent}, there is no agent--graph
pair $(\Agent, \Graph)$ admitting a realised distinguishability
act of residue $< \floor$, where $\floor = \min\{w(e) : e \in \Edg,
w(e) > 0\}$ is the graph-level floor of \cref{thm:floor}.
\end{theorem}

\begin{proof}
Suppose for contradiction such a pair exists, with a realised act
$u \not\sim v$ of residue $\Res(u, v) < \floor$. The residue is the
sum of edge measures along the agent's path. Each edge of the path
has measure $0$ or $\geq \floor$ (by the definition of $\floor$ as
the minimum positive edge measure in $\Graph$). Hence the sum is
$0$ or $\geq \floor$.

If the sum is $0$, every edge on the path has measure $0$. The
agent's descriptive scheme has no measure differential to record
the distinction (\cref{F:fin_descr}, see the realised-path argument
in the proof of \cref{thm:floor}), so the act is not realised in
the agent's record. The assumed realised act does not exist.

If the sum is $\geq \floor$, the residue is $\geq \floor$,
contradicting the assumption $\Res(u, v) < \floor$.

Either case contradicts the assumption. $\qed$
\end{proof}

\begin{corollary}[The main theorem is irrefutable from within the
premises]
\label{cor:irrefutable}
\Cref{thm:opstructure} cannot be refuted by an example satisfying
\cref{ax:graph} and \cref{ax:agent}. Refutation requires either
denying one of the premises or finding an error in a constituent
lemma; no example produced under the premises is a counterexample.
\end{corollary}

\begin{proof}
Refutation by example requires a $(\Agent, \Graph)$ satisfying the
premises in which the theorem fails. By the construction of
$\OpStr_\Agent$ from $(\Agent, \Graph)$ in the theorem's proof, and
the zero-floor impossibility (\cref{thm:zero_floor}), no such
example exists. $\qed$
\end{proof}

\subsection{Honest scope}

\begin{remark}[The argument is silent outside its premises]
\label{rem:silent}
We make no claim about agents that fail \cref{ax:agent} (putative
infinite agents, oracles, perfect observers) or systems that fail
\cref{ax:graph} (continuous fields without finite discretisation,
infinite combinatorial structures). The argument has nothing to
say about those cases. This silence is intended: the framework is
deliberately narrow, sharp where the premises hold, mute elsewhere.
This is the right shape for a conditional architectural claim. It
allows the result to be definite without overreaching.
\end{remark}

% =====================================================================
\section{Implications}
\label{sec:implications}
% =====================================================================

We draw three implications of the main theorem for engineering
practice. Each is a direct consequence of \cref{thm:opstructure} or
its corollaries.

\subsection{Software architecture under the conditional}

The conjunction of \cref{ax:graph} and \cref{ax:agent} covers all
realised software systems: any program with finite memory and
finite execution is a finite agent, and any data structure it
manipulates is a finite graph. By \cref{thm:opstructure}, the
program's operational structure is forced by its data structure and
its agency: it is not chosen.

In practice this means: arguments about which operations a program
``should'' admit reduce to arguments about $(\Agent, \Graph)$.
Two architectures disagreeing about the operational interface are
disagreeing about one of: which finite data structure the system is
modelled as; which finite agent the program is. Once those are
agreed, the operations are determined.

\subsection{Database query algebras}

A finite database is a finite graph in the sense of
\cref{def:graph}: the schema gives the vertices (tables, attributes,
values), the foreign-key relations give the edges, and the
cardinalities give the edge measures. A finite query engine is a
finite agent in the sense of \cref{def:agent}: bounded internal
state, finite query plan space, finite cost budget per query.

By \cref{thm:opstructure}, the query algebra is forced. The choice
between relational algebra, calculus, and CRUD is the choice of
which subset of the forced algebra to expose; it is not the choice
of which algebra to invent. The minimal operational floor is the
smallest positive edge weight in the schema graph (typically a unit
cardinality), and every realisable query has residue $\geq$ this
floor.

\subsection{Knowledge graphs and reasoning}

A finite knowledge graph---ontology, RDF triple store, labelled
property graph---is a finite graph in the sense of \cref{def:graph}.
A finite reasoner---SPARQL engine, description-logic prover, graph
neural network---is a finite agent in the sense of \cref{def:agent}.
By \cref{thm:opstructure}, the reasoner's operational structure is
forced by the graph it queries and the resources it has.

The implication is that the apparent diversity of reasoning systems
collapses, under the conditional, to relabelling. Two reasoners
with the same $(\Agent, \Graph)$ admit the same operations up to
naming; their differences are presentational, not structural.
Cross-system calibration (translation of one query language into
another) is guaranteed to exist whenever $(\Agent, \Graph)$ is
shared.

% =====================================================================
\section{Conclusion}
\label{sec:conclusion}
% =====================================================================

\subsection{What we have established}

We have shown that under two finiteness premises---finite engaging
agent (\cref{ax:agent}), finite engaged graph-structured system
(\cref{ax:graph})---the system necessarily admits an operational
structure $\OpStr_\Agent$ that is forced by the premises and unique
given the pair $(\Agent, \Graph)$ together with the agent's
path-choice policy. The structure is a small category with partial
composition, a residue function with a strictly positive floor on
non-identity morphisms, a residue preorder that descends to a
partial order on the residue quotient, a budget-bounded
realisability marker, and a four-type classification (algebraic,
dynamical, computational, internal) that exhausts the realisable
operations.

The result is a conditional architectural claim. It does not
address what reality is in itself; it addresses what structure
follows when the two conditions hold. Where they hold---which is
the case in essentially all realised engineering systems---the
operational structure is determined.

\subsection{What we have not claimed}

We have not claimed that:
\begin{enumerate}[nosep,leftmargin=*]
\item Reality is fundamentally finite. The premises restrict the
agent and the system the agent engages, not the world.
\item All systems are graphs. The premises identify the conditions
under which the result holds; systems failing the conditions are
not addressed.
\item Operational structures are the only structures of interest.
Other structures (continuous dynamics, categorical limits, infinite
processes) may coexist with operational structure; we have not
spoken to them.
\end{enumerate}

\subsection{The narrowness of the result is the point}

A conditional result is sharp because it is narrow. The framework
established here speaks definitively about a specific class of
situations and is silent about others. This is exactly the right
shape for an architectural claim: scope is a feature, and the scope
of \cref{thm:opstructure} is the scope of finite agency engaging
finite graph-structured systems. Within that scope, the operational
structure is forced. Outside it, the framework declines to address
the case.

\subsection{What this enables}

The main practical consequence is that engineering decisions over
$(\Agent, \Graph)$ become discovery tasks rather than invention
tasks. To design a finite system's operational structure is to
identify the unique $\OpStr_\Agent$ that the agent--graph pair
determines. To compare two designs is to ask whether they realise
the same $(\Agent, \Graph)$. To extend a system is to ask how the
extension changes $(\Agent, \Graph)$ and re-derive the structure.

The framework converts what is currently treated as a matter of
architectural taste into a matter of formal derivation. Whether
this is a useful conversion in practice depends on whether the cost
of formalising $(\Agent, \Graph)$ for a given system is less than
the cost of debating its operational structure informally. We
expect this to be true at scale---for large systems with many
stakeholders and long lifetimes---and false at small scale, where
the informal debate is cheap.

\subsection{Closing}

The structure of operations on a finite graph engaged by a finite
agent is, in its uniquely-determined part, not a design problem. It
is a derivation from the two finiteness premises. The derivation
produces an answer unique up to labelling and path-choice: a small
category with a positive floor, a residue preorder, partial
composition under a budget constraint, and a four-type
classification of effects. That is what the conditional
architecture says, and that is all it says.

% =====================================================================
\begin{thebibliography}{9}

\bibitem[Bollob\'as(1998)]{bollobas1998} B\'ela Bollob\'as.
\emph{Modern Graph Theory}.
Graduate Texts in Mathematics 184, Springer, 1998.

\bibitem[Davey and Priestley(2002)]{daveypriestley2002}
B.~A. Davey and H.~A. Priestley.
\emph{Introduction to Lattices and Order}, 2nd edition.
Cambridge University Press, 2002.

\bibitem[Diestel(2017)]{diestel2017} Reinhard Diestel.
\emph{Graph Theory}, 5th edition.
Graduate Texts in Mathematics 173, Springer, 2017.

\bibitem[Halmos(1950)]{halmos1950} Paul R. Halmos.
\emph{Measure Theory}.
Van Nostrand, 1950.

\bibitem[Hopcroft, Motwani and Ullman(2007)]{hopcroft2007}
John E. Hopcroft, Rajeev Motwani, and Jeffrey D. Ullman.
\emph{Introduction to Automata Theory, Languages, and Computation},
3rd edition. Addison-Wesley, 2007.

\bibitem[Howie(1995)]{howie1995} John M. Howie.
\emph{Fundamentals of Semigroup Theory}.
London Mathematical Society Monographs 12, Oxford University Press,
1995.

\bibitem[Mac Lane(1998)]{maclane1998} Saunders Mac Lane.
\emph{Categories for the Working Mathematician}, 2nd edition.
Graduate Texts in Mathematics 5, Springer, 1998.

\end{thebibliography}

\end{document}
